% ===== tr-02 块2：第 3 章 pp.19--21（原文页 19--21）=====
% [原文 p.19]
（一个例子见 [44]。）所以，我们将只考察特殊拉格朗日环面的模空间。

% [Tutor Note: 原句 "…nearby special Lagrangian submanifolds may intersect." 跨 p.18/p.19，按「开始者写完」由 tr02e 写完整句；本文件从 p.19 的第二句 "So instead, we will just look at…" 开始，不重复、不遗漏。]

假设给定 $t_i\to0$，且对充分接近 0 的 $t_i$，存在一个特殊拉格朗日 $T^n$，其同调类在单值变换下不变；更具体地说，它生成单值权重滤过（monodromy weight filtration）的空间 $W_0$（我们正是在这里期待找到与极大幂零退化相伴随的特殊拉格朗日纤维化的纤维）。设 $B_{0,i}$ 是这个环面的形变模空间；$B_{0,i}$ 的每个点对应于 $\mathcal{X}_{t_i}$ 中一个光滑的特殊拉格朗日环面。于是这个流形带上了第 2 节定义的 McLean 度量与仿射结构。然后我们可以紧化 $B_{0,i}\subseteq B_i$（大概是通过取 $B_{0,i}$ 在特殊拉格朗日流（current）空间中的闭包；细节在此并不重要）。这给出一列度量空间 $(B_i,d_i)$，其度量 $d_i$ 由 McLean 度量诱导。若把 McLean 度量规范化，使 $B_i$ 的直径与 $i$ 无关地保持常数，则我们可以希望 $(B_i,d_i)$ 收敛到一个紧度量空间 $(B_\infty,d_\infty)$。于是 SYZ 的极限形式如下：

\begin{conjecture}{3.5.5}
若 $(\mathcal{X}_i,g_i)$ 收敛到 $(X_\infty,g_\infty)$，且对大的 $i$，$(B_i,d_i)$ 非空并收敛到 $(B_\infty,d_\infty)$，则 $B_\infty$ 与 $X_\infty$ 在相差一个缩放的意义下等距。此外，存在子空间 $B_{\infty,0}\subseteq B_\infty$，使得 $\Gamma:=B_\infty\setminus B_{\infty,0}$ 在 $B_\infty$ 中是 Hausdorff 余维 2 的，并且 $B_{\infty,0}$ 是一个 Monge--Ampère 流形，其 Monge--Ampère 度量在 $B_{\infty,0}$ 上诱导出 $d_\infty$。
\end{conjecture}

% [Tutor Note: 猜想 5.5 中 $B_\infty$ 与 $X_\infty$ 的关系是 "isometric up to scaling"（相差缩放意义下等距），不是严格等距；这是 SYZ 极限图景的核心断言------底空间即 Gromov--Hausdorff 极限。注意 $\Gamma$ 是相对于 $B_\infty$ 的 Hausdorff 余维 2 集合，而不是 $B_\infty$ 的拓扑边界。]

本质上，这说的是：当趋近极大幂零退化时，我们期待在 $\mathcal{X}_{t_i}$ 越来越大的子集上拥有特殊拉格朗日纤维化。此外，在极限中，Joyce 所提示的余维 1 判别轨迹收敛为余维 2 的判别轨迹，而 $B_{0,i}$ 上（不一定为 Monge--Ampère 的，见 [44]）Hessian 度量收敛为 Monge--Ampère 度量。

我在这里想要说明的主要之点是：SYZ 猜想很可能只是「近似地」正确，而人们必须考察极限，才有证明任何东西的希望。另一方面，上述猜想看起来很可能用目前所理解的技术就可以达到，尽管还需要大量额外工作；若在今后几年内看到它被证明，我不会感到惊讶。

我们如何用 SYZ 猜想的这一修正版本做镜像对称？本质上，我们会遵循以下步骤：

\begin{enumerate}
\item 从 Calabi--Yau 流形的一个极大幂零退化 $\mathcal{X}\to D$ 出发，并选取一个极化。这给出对每个 $t\in D\setminus0$ 的一个 Kähler 类 $[\omega_t]\in H^2(\mathcal{X}_t,\R)$，它由 $\omega_t$ 代表，而 $\omega_t$ 是某个 Ricci 平坦度量 $g_t$ 的 Kähler 形式。
\item 识别序列 $(\mathcal{X}_{t_i},r_ig_{t_i})$ 的 Gromov--Hausdorff 极限，其中 $t_i\to0$，而 $r_i$ 是使 $\mathcal{X}_{t_i}$ 的直径保持不变的缩放因子。若上述猜想成立，该极限将是一个带奇点的仿射流形 $B$ 连同 Monge--Ampère 度量。
\item 作 Legendre 变换，得到一个新的带奇点仿射流形 $\check B$，但度量不变。
\item 尝试对小 $\epsilon>0$ 构造 $X_\epsilon(\check B_0)$ 的紧化，以得到一个复流形 $X_\epsilon(\check B)$。这将就是镜像流形。
\end{enumerate}

实际上，我们需要对最后一步再多作些说明。问题在于：虽然我们期待在一般情况下可以构造
% [原文 p.20]
辛流形 $\check X(B_0)$ 的辛紧化（从而把镜像作为辛流形得到），但我们并不期待能把 $X_\epsilon(\check B_0)$ 紧化为复流形。相反，预期是在它能被紧化之前，必须先对 $X_\epsilon(\check B_0)$ 作一个小形变。此外，这个小形变在镜像对称中至关重要：正是这个小形变提供了 \emph{$B$-模型瞬子修正}。

% [Tutor Note: 第 (4) 步与后文两处 $X_\epsilon(\check B_0)$、$X_\epsilon(\check B)$ 的 check 记号极易漏抄：$(4)$ 的结论是"紧化 $X_\epsilon(\check B_0)$ 得到复流形 $X_\epsilon(\check B)$"；文字层会把这些 check 丢掉，须对照 p.19 末与 p.20 页图（700 dpi 裁图复核）。]

因为最后这个问题很重要，它已经被若干作者研究过：Fukaya 在 [11] 中用启发式想法直接研究了这一问题；而 Kontsevich 与 Soibelman [39] 修改了从仿射流形过渡到复流形的问题，改为产生一个非阿基米德空间。我们将在本文后面讨论我自己与 Siebert 的工作时回到这些问题，那部分工作正部分地由同一问题所推动。由于最后这一项如此重要，让我们给它起个名字：

\begin{question}{3.5.6}[重建问题（版本 1）]
给定一个带奇点的热带仿射流形 $B$，构造一个复流形 $X_\epsilon(B)$，它是 $X_\epsilon(B_0)$ 的某个小形变的紧化。
\end{question}

我不愿在此进一步停留于 SYZ 猜想的这个版本，因为它主要属于分析、微分几何以及 Ricci 平坦度量行为的领域，对于我们理解传统镜像对称中更微妙的方面（例如瞬子修正究竟如何产生）帮助不大。所以我们转而解释 SYZ 的极限形式如何启发了 SYZ 的一个更代数几何化的形式；后者反过来避开了所有分析问题，并为理解镜像对称的根本奥秘带来了很大的希望。

% [Tutor Note: 问题 5.6 中的 "$X_\epsilon(B)$" 与 "$X_\epsilon(B_0)$" 都不带 check（与第 (4) 步的 $X_\epsilon(\check B)$ 不同）；原文在同一段里混用带 check 与不带 check 的记号，译文逐一照抄，不做统一。]

\section{Gromov--Hausdorff 极限、代数退化与镜像对称}

我们现在有两种极限概念：一方面是熟悉的代数几何概念，即圆盘上的平坦退化族 $\mathcal{X}\to D$；另一方面是 Gromov--Hausdorff 极限。Kontsevich 对这两者之间的联系有一个重要的洞见（见 [38]）。本节将粗略说明这一联系如何以及为何成立。

非常粗略地说，当 $t_i\to0$ 时，Gromov--Hausdorff 极限 $(\mathcal{X}_{t_i},g_{t_i})$ 在拓扑上应当与奇异纤维 $\mathcal{X}_0$ 的对偶相交复形重合。更确切地说，在一个相对简单的情形中，假设 $f:\mathcal{X}\to D$ 是相对极小（Mori 意义下）且为正规交叉的，并且 $\mathcal{X}_0$ 的不可约分支为 $X_1,\dots,X_m$。$\mathcal{X}_0$ 的对偶相交复形是以 $v_1,\dots,v_m$ 为顶点的单纯复形，它包含单形 $\langle v_{i_0},\dots,v_{i_p}\rangle$ 当且仅当 $X_{i_0}\cap\cdots\cap X_{i_p}\ne\varnothing$。

让我们粗略解释为什么应当如此：先看一族标准的退化椭圆曲线，其周期为 1 与 $\frac{n}{2\pi i}\log t$，其中 $n$ 是正整数。这样一族曲线位于去心圆盘上，通过在原点处添加一个 $I_n$ 纤维（$n$ 条有理曲线组成的一个循环）延拓为圆盘上的一族。

% [Tutor Note: $I_n$ 是 Kodaira 的 $I_n$ 型奇异纤维：由 $n$ 条有理曲线围成一个 $n$-边形循环；p.15 处理过的 $I_1$ 型就是被捏住的环面（$n=1$）。这里 $n$ 既用作周期公式里的整数，又用作 p.21 中不可约分支的个数，注意区分。]

取 $t_i\to0$ 且 $t_i$ 为实数、正数，得到一列形式为 $X_{\epsilon_i}(B)$ 的椭圆曲线，其中 $B=\R/n\Z$，$\epsilon_i=-\frac{2\pi}{\ln t_i}$。此外，$X_{\epsilon_i}(B)$ 上的度量经适当缩放后来自 $B$ 上的常值 Hessian 度量。所以我们想解释 $B$ 与奇异纤维附近的几何如何相关联。为此，
% [原文 p.21]
设 $X_1,\dots,X_n$ 是 $\mathcal{X}_0$ 的不可约分支；它们全都是 $\PP^1$。设 $P_1,\dots,P_n$ 是 $\mathcal{X}_0$ 的奇异点。

我们考虑 $\mathcal{X}$ 中两类开集。对第一类，在 $X_i$ 上取坐标 $z$，使得 $P_i$ 由 $z=0$ 给出、$P_{i+1}$ 由 $z=\infty$ 给出。设 $U_i\subseteq D_i$ 为开集 $\{z\mid\delta\le|z|\le1/\delta\}$，其中 $\delta$ 是某个固定的小量。于是可以找到 $U_i$ 在 $\mathcal{X}$ 中的邻域 $\tilde U_i$，使得 $\tilde U_i$ 双全纯于 $U_i\times D_\rho$（$\rho>0$ 充分小，$D_\rho$ 是 $\C$ 中半径为 $\rho$ 的圆盘），并且 $f|_{\tilde U_i}$ 是到 $D_\rho$ 的投影。

另一方面，每个 $P_i$ 在 $\mathcal{X}$ 中有一个邻域 $\tilde V_i$，双全纯于多圆盘 $\{(z_1,z_2)\in\C^2\mid |z_1|\le\delta',\,|z_2|\le\delta'\}$，在其上 $f$ 取形式 $z_1z_2$。

若 $\delta$ 与 $\delta'$ 选取得当，则对充分接近 0 的 $t$，
\[
\{\tilde V_i\cap\mathcal{X}_\epsilon \mid 1\le i\le n\}\cup\{\tilde U_i\cap\mathcal{X}_t \mid 1\le i\le n\}
\]
构成 $\mathcal{X}_t$ 的一个开覆盖。现在该开覆盖中的每个集合都可以写成 $X_\epsilon(U)$，其中 $U$ 是一维（非紧）仿射流形，$\epsilon=-2\pi/\ln|t|$。若 $U$ 是开区间 $(a,b)\subseteq\R$，则 $X_\epsilon(U)$ 双全纯于圆环
\[
\{z\in\C \mid e^{-2\pi b/\epsilon}\le|z|\le e^{-2\pi a/\epsilon}\},
\]
因为 $q=e^{2\pi i(x+iy)/\epsilon}$ 是 $X_\epsilon((a,b))$ 上的全纯坐标。因此
\[
\tilde U_i\cap\mathcal{X}_t\cong X_\epsilon\Big(\Big(\frac{\epsilon\ln\delta}{2\pi},-\frac{\epsilon\ln\delta}{2\pi}\Big)\Big)
\]
其中 $\epsilon=-2\pi/\ln|t|$。当 $t\to0$ 时，区间 $(\epsilon\ln\delta/2\pi,-\epsilon\ln\delta/2\pi)$ 收缩成一个点。所以在这种看法下，当 $t\to0$ 时，$\tilde U_i\cap\mathcal{X}_t$ 是 $\mathcal{X}_t$ 中越来越小的开子集。这一论证提示：每个不可约分支都应当与 $B$ 上的一个点相关联。

现在考察 $\tilde V_i\cap\mathcal{X}_t$。它是
\begin{align*}
\{(z_1,z_2)\in\C^2 \mid |z_1|,|z_2|<\delta',\ z_1z_2=t\}&\cong \{z\in\C \mid |t|/\delta'\le|z|\le\delta'\}\\
&\cong X_\epsilon\Big(-\frac{\epsilon}{2\pi}\ln\delta',\ \frac{\epsilon}{2\pi}(\ln\delta'-\ln|t|)\Big)
\end{align*}
其中 $\epsilon=-2\pi/\ln|t|$。当 $t\to0$ 时，这个区间趋近单位区间 $(0,1)$。于是对小的 $t$，开集 $\tilde V_i\cap\mathcal{X}_t$ 最终成为 $\mathcal{X}_t$ 的很大一部分。最终我们看到：对小的 $t$，$\mathcal{X}_t$ 是若干形如 $X_\epsilon((i+\epsilon',i+1-\epsilon'))$（即 $\tilde V_i\cap\mathcal{X}_t$）与形如 $X_\epsilon((i-\epsilon',i+\epsilon''))$（即 $\tilde U_i\cap\mathcal{X}_t$）的开集之并，其中 $\epsilon',\epsilon''$ 充分小。这些开集应当（至少近似地）粘合起来，给出 $X_\epsilon(B)$。于是我们看到：$\mathcal{X}_0$ 的不可约分支似乎与 $B$ 上的点重合，而分支之间的相交则与直线重合。这样我们就看到对偶相交复形浮现出来。

% [Tutor Note: 符号核对：$\epsilon=-2\pi/\ln|t|>0$（因为 $|t|<1$ 时 $\ln|t|<0$），而 $\ln\delta<0$，故 $\epsilon\ln\delta/2\pi<0<-\epsilon\ln\delta/2\pi$，区间 $(\epsilon\ln\delta/2\pi,-\epsilon\ln\delta/2\pi)$ 的两端确实满足左小右大；当 $t\to0$ 时 $\epsilon\to0^+$，该区间收缩成一点，于是 $\tilde U_i\cap\mathcal{X}_t$ 在 $\mathcal{X}_t$ 中越来越小。]

% [Tutor Note: 末段「每个不可约分支 ↔ B 上的点，分支之间的相交 ↔ 直线」是本节的核心结论：$\mathcal{X}_0$ 的对偶相交复形的顶点数等于不可约分支数，其边数等于相交数，于是 Gromov--Hausdorff 极限的骨架正是 $B$。]

让我们在下一节开始严格结果之前再作一个观察。更一般地，假设我们有 Calabi--Yau 流形的一个 Gorenstein 环面交叉退化 $f:\mathcal{X}\to D$。这意指每一点 $x\in\mathcal{X}$ 都有一个同构于某个仿射 Gorenstein 环面簇（即典范类是 Cartier 除子）中开集的邻域，而 $f$ 局部地由某个单项式给出，该单项式在每个余维 1 的环面分层上恰好消失到 1 阶。这是正规交叉概念的推广，见 [55]。非常粗略地说，上述论证提示：中心纤维的每个不可约分支都将对应于 Gromov--Hausdorff 极限的一个点。下面的习题说明了当 $x\in\mathcal{X}_0$ 是 $\mathcal{X}_0$ 中的零维分层时，应当期待它对 $B$ 作出何种贡献。

% [Tutor Note: 本块（b04, pp.11--21）到此结束：p.21 末句以句号收尾，p.22 首行为 "Exercise 6.1."（归 tr-03 b05）。已核对 p.22 split 文本，无跨界句需要本块续写。]
